\section{Analytic evaluation of cut master integrals}
\label{sec:analytic-masters}

IBP reduction determines the rational coefficients but does not evaluate the
masters.  For each master we retain the physical cut momenta, their energy
orientations, the ordinary propagators, the physical chamber, and the
normalization of the integration measure.  The required output is a
meromorphic analytic function of $\eps$ and the invariants.

\subsection{Physical parameterizations and branches}

Momentum-conservation delta functions are eliminated exactly, after which a
Lorentz frame or a Feynman-parameter representation may be chosen for the
remaining integrations.  Every change of variables is accompanied by its
Jacobian and by the image of the physical domain.  Before any Laurent
expansion, each noninteger power is written as
\begin{equation}
 X^\alpha=\exp[\alpha\log X]
 \label{eq:real-power-definition}
\end{equation}
with the sign of $X$ fixed in the stated chamber.  If a powered base crosses
zero inside the contour, the integral is divided into regions or treated with
its explicit causal boundary value.  Algebraic replacements such as
$(-X)^\alpha\to X^\alpha$ are never made without this analysis.

\subsection{Direct integration and differential equations}

When the reduced parameter density is linearly reducible, its Laurent
coefficients are integrated exactly with sector decomposition followed by
hyperlogarithmic integration.  \SubTropica\ is used for the analytic sector
expansion and \HyperInt-compatible functions for the resulting iterated
integrals~\cite{Giroux:2026tgd,Panzer:2014caa}.  This route returns analytic
coefficients, not values at fixed kinematics.

When direct integration is not efficient, a vector of masters
$\boldsymbol M(\boldsymbol x,\eps)$ is differentiated with respect to the
physical invariants and reduced by the same IBP rules:
\begin{equation}
 \dd\boldsymbol M=A(\boldsymbol x,\eps)\boldsymbol M.
 \label{eq:master-differential-system}
\end{equation}
The exact integrability condition
\begin{equation}
 \dd A-A\wedge A=0
 \label{eq:flatness-condition}
\end{equation}
is checked before solving the system.  Rational transformations are then used
to reach Fuchsian or canonical form when possible
\cite{Henn:2013pwa,Meyer:2017joq}.  The solution is fixed by analytic boundary
data obtained from a simpler cut phase space, a regular kinematic point, or a
controlled expansion by regions.  A numerical boundary value may check this
derivation but does not replace it.

\subsection{Four distinct relations among masters}

Several reductions of the boundary workload are useful, but they are not the
same statement.  Let $I_i(\boldsymbol x,\eps)$ denote powered cut integrals.

First, $I_i\sim_{\mathrm{pow}}I_j$ when an invertible momentum change and an
allowed external relabeling map the complete measure, contour, oriented cuts,
ordinary denominators, and every propagator power of $I_i$ into those of
$I_j$.  This is equality of the integrals after the corresponding kinematic
substitution.

Second, $I_i\sim_{\mathrm{den}}I_j$ when the same kind of map identifies the
sets of positive-power denominator slots while powers are ignored.  This
classifies denominator geometry but does not prove equality of powered
integrals.

Third, a set $S$ is closed under differentiation in variables
$\boldsymbol x$ when
\begin{equation}
 \partial_{x_r}I_i\in
 \operatorname{span}_{\mathbb Q(\boldsymbol x,\eps)}S
 \qquad
 \text{for every }I_i\in S\text{ and every }x_r.
 \label{eq:de-closure-definition}
\end{equation}
This determines which masters must be solved together.  Differential coupling
does not imply that their boundary constants are equal.

Finally, two scalar boundary integrals represent the same boundary period
when an invertible change of integration variables maps the contour, measure,
integrand, normalization, and branch of one into the other.  We denote this by
$B_i\sim_{\partial}B_j$.  This is the relation that reduces the number of
independent analytic boundary integrations.  Section~\ref{sec:nlo-example}
shows explicitly that six NLO masters occupy two maximal nontrivial
denominator classes but require only one non-elementary boundary period once
the cut bubble is known.

\subsection{Independent numerical comparison}

For a physical point away from singular surfaces, the analytic Laurent
coefficients are compared with \AMFlow\ values generated from the same family
and normalization~\cite{Liu:2020kpc}.  Agreement is required coefficient by
coefficient to a declared precision.  This tests the analytic expression and
its branch choice; it does not determine the function or supply missing
boundary constants.
